% Generated draft: requires mathematical and source review.
\documentclass[11pt]{article}
\usepackage[margin=25mm]{geometry}
\usepackage{fontspec}
\setmainfont{Latin Modern Roman}
\usepackage{amsmath,amssymb,mathtools,unicode-math}
\setmathfont{Latin Modern Math}
\usepackage{enumitem,needspace,xurl,hyperref}
\setlength{\parindent}{0pt}
\setlength{\parskip}{5pt}
\setlength{\emergencystretch}{3em}
\newcommand{\N}{\mathbb{N}}
\newcommand{\Z}{\mathbb{Z}}
\newcommand{\Q}{\mathbb{Q}}
\newcommand{\R}{\mathbb{R}}
\newcommand{\C}{\mathbb{C}}
\newcommand{\F}{\mathbb{F}}
\newcommand{\A}{\mathbb{A}}
\newcommand{\PP}{\mathbb{P}}
\newcommand{\E}{\mathbb{E}}
\newcommand{\eps}{\varepsilon}
\newcommand{\dd}{\,\mathrm{d}}
\begin{document}
\section*{UnsolvedMath ID: 10600003}
\textbf{Source problem code:} catalogue entry
\subsection{Definitions and mathematical statement}
Let \(\alpha\) be any ordinal, viewed as the ordered set of ordinals smaller than \(\alpha\). The source's flat hierarchy uses virtual diagrams whose flat crossings are labeled by members of \(\alpha\); virtual crossings are distinct from these labeled flat crossings. A flat crossing carries no over/under choice.

The supplied third-Reidemeister rule is this: a strand passing through two crossings both labeled \(a\) may slide across a crossing labeled \(b\) only when
\[
 a>b,\qquad a,b\in\alpha.
\]
The inequality is strict. The source says this produces arbitrarily many types of flat crossings and generalizes ordinary flat virtual diagram theory.

Problem-list item 3 asks to classify the diagrams of this hierarchy, for arbitrary \(\alpha\). To make that an equivalence-classification problem, the full generating moves must be specified. The record gives no label rules for first and second Reidemeister moves, no complete list of third moves, and no explicit account of ordinary same-label moves or virtual detour moves. In particular the displayed permission excludes \(a=b\), so its compatibility with a purported single-type ordinary flat theory must be checked rather than repaired by replacing \(>\) with \(\ge\). The source also does not specify closed versus long diagrams, component counts, orientations, or the desired form of classification. Thus the ordinal domain and given local rule are precise, but the whole equivalence relation remains an unresolved source dependency.
\subsection{Short English statement}
Classify flat virtual diagrams whose flat crossings carry ordinal labels and whose specified third move is allowed only when the repeated label is strictly larger than the crossed label. The other move rules and the intended classification format are missing. The concept is attributed to Kauffman as unpublished, with an empty first-step citation.
\subsection{Sources}
\begin{itemize}
\item Local UnsolvedMath registry, permanent ID \(10600003\), exact source code \texttt{catalogue entry}, ``Virtual-knot problem 3 --- The Flat Hierarchy.'' The statement attributes the concept to Kauffman, unpublished; no proposed year is supplied.
\item \emph{Unsolved Problems in Virtual Knot Theory and Combinatorial Knot Theory}, revised survey, dedicated problem-list item 3: \url{https://arxiv.org/abs/1409.2823}. The supplied background describes 60 top-level enumerated problems with compound items retained intact.
\item Source-TeX extraction, accessed 2026-07-29. The sentence ``A first step in working with the flat hierarchy can be found in .'' has an empty citation, which is not reconstructed.
\item Source evidence NEEDS\_REVIEW presents an open problem and requires a current primary-literature audit. Supplied OPEN-TRIAGE assessment reports no located literature on this specific ordinal-labeled hierarchy and describes it as apparently open and unstudied. This is the supplied bounded-search assessment, not an independent verification.
\item Dataset attribution: Ulam AI / UnsolvedMath Contributors, CC BY 4.0. Source rights NEEDS\_REVIEW: publicly accessible author/publisher copy, redistribution terms requiring review.
\end{itemize}
\end{document}
